% !TeX root = ../navier-stokes-manuscript.tex
% ============================================================
% 2.7 Compatibility and the Terminal Quantifier
% ============================================================

\subsection{Compatibility and the Terminal Quantifier}\label{sub:CompatibilityAndTheTerminalQuantifier}

Finite rectangles can recover one intersection only if cutting the mixed jet at different depths never changes a coordinate they share.
The first nontrivial comparison can be seen directly.
Take
\[
  \cR_{2,3}[u_0;T].
\]
Its upper velocity level contains
\[
  V_0,V_1,V_2
  \in
  L^2(0,T;H^4),
\]
and its lower velocity level contains
\[
  V_1,V_2,V_3
  \in
  L^2(0,T;H^2).
\]
It also retains \(Q_0,Q_1,Q_2\) with the corresponding momentum, divergence, and pressure equations.

Now keep only temporal orders through \(n=1\) and lower the middle trace height from \(H^3\) to \(H^2\).
The continuous embeddings
\[
  H^4\hookrightarrow H^3,
  \qquad
  H^2\hookrightarrow H^1
\]
give
\[
  V_0,V_1
  \in
  L^2(0,T;H^3),
  \qquad
  V_1,V_2
  \in
  L^2(0,T;H^1).
\]
Keeping \(Q_0,Q_1\) and their equations leaves exactly the coordinates of
\[
  \cR_{1,2}[u_0;T].
\]
The surviving \(V_1\) is the same \(\partial_tu\), the surviving \(Q_1\) is the same \(\partial_tp\), and every surviving spatial derivative is the same derivative of those fields.
The smaller cut changes only the retained depth, not any retained derivative.

The general restriction has independent temporal, spatial, and time cutoffs.
Let
\[
  N'\le N,
  \qquad
  M'\le M,
  \qquad
  0<T'\le T<T_*.
\]
Define
\[
  \rho_{N',M';T'}^{N,M;T}:
  \mathscr X_{N,M}((0,T))
  \longrightarrow
  \mathscr X_{N',M'}((0,T'))
\]
by keeping \(V_0,\ldots,V_{N'+1}\), keeping \(Q_0,\ldots,Q_{N'}\), restricting every retained field to \((0,T')\), and reading it through the embeddings from the \(M\)-levels to the \(M'\)-levels.
The smaller value
\[
  \mathfrak R_{N',M'}(u;T')
\]
is recomputed from those restricted fields with its own indices and interval.
It is not asserted to equal the numerical value of \(\mathfrak R_{N,M}(u;T)\).
The embeddings and time restriction assert only that the recomputed smaller norm is finite whenever the larger one is finite.
For the datum-generated records, this map satisfies
\[
  \rho_{N',M';T'}^{N,M;T}
  \bigl(\cR_{N,M}[u_0;T]\bigr)
  =
  \cR_{N',M'}[u_0;T'].
\]

For every shared temporal and spatial index,
\[
  J_{n,\alpha}
  =
  \partial_t^n\partial_x^\alpha u,
  \qquad
  \Pi_{n,\alpha}
  =
  \partial_t^n\partial_x^\alpha p
\]
remain the same distributional derivatives.
The retained rows continue to satisfy the corresponding pressure-complete recurrence because they are restrictions of the same equation-generated jet.
Compatibility is thus structural.
It is not a later agreement imposed on rectangles that were generated independently.

Restriction in stages agrees with direct restriction.
If
\[
  N''\le N'\le N,
  \qquad
  M''\le M'\le M,
  \qquad
  0<T''\le T'\le T<T_*,
\]
then
\[
  \rho_{N'',M'';T''}^{N',M';T'}
  \circ
  \rho_{N',M';T'}^{N,M;T}
  =
  \rho_{N'',M'';T''}^{N,M;T}.
\]
Every route to a shared lower coordinate returns the same derivative of the same \((u,p)\).
The family of finite cuts can consequently represent one velocity intersection with its attached pressure rows.

Structural agreement does not yet decide whether a fixed rectangle survives the approach to \(T_*\).
Fix \(T<T_*\).
Classical smoothness on the compact interval \([0,T]\) gives, for every finite \(n\) and \(q\),
\[
  \sup_{0\le t\le T}
  \norm{\partial_t^nu(t)}_{H^q}
  <
  \infty.
\]
Hence
\[
  \norm{\partial_t^nu}_{L^2(0,T;H^q)}^2
  \le
  T
  \sup_{0\le t\le T}
  \norm{\partial_t^nu(t)}_{H^q}^2
  <
  \infty.
\]
Every finite rectangle is finite after a strict right endpoint has been chosen.
The quantifiers are
\[
  \forall\,0<T<T_*
  \ \forall\,N,M\in\Nzero
  \ \exists\,C_{T,N,M}<\infty:
  \qquad
  \mathfrak R_{N,M}(u;T)
  \le
  C_{T,N,M}.
\]
This statement permits \(C_{T,N,M}\) to diverge as \(T\uparrow T_*\).
Strict-cutoff finiteness does not supply terminal control.

Whole-terminal control reverses the decisive part of the order.
First fix one finite coordinate pair \((N,M)\), and then require one constant to survive every strict cutoff:
\[
  \forall\,N,M\in\Nzero
  \ \exists\,C_{N,M}(u_0,\nu,T_*)<\infty
  \ \forall\,0<T<T_*:
  \qquad
  \mathfrak R_{N,M}(u;T)
  \le
  C_{N,M}(u_0,\nu,T_*).
\]
Equivalently,
\[
  \sup_{0<T<T_*}
  \mathfrak R_{N,M}(u;T)
  \le
  C_{N,M}(u_0,\nu,T_*)
  <
  \infty
  \qquad
  \text{for every finite pair }(N,M).
\]
The constant may depend on \(N\) and \(M\).
Uniformity concerns the cutoff \(T\uparrow T_*\) for one fixed rectangle, not derivative order across all rectangles.

\Cref{thm:WholeTerminalCompatibleRectangleEstimate} supplies this whole-terminal bound for every finite pressure-complete rectangle generated by the datum \(u_0\).
Every rectangle is already a cut through the original equation-generated jet, so the bounds carry that compatible family to the terminal cutoff without inventing its agreement.
Allowing all finite pairs gives the complete velocity intersection first; its fixed \(H^{1/2}\) and \(H^{3/2}\) rows then recover the critical quantities.
